% p.84 — طريقۂ تعليم (Urdu: the author's note on teaching method)
% Layout: centred large heading with a rule beneath it; one Urdu paragraph (first line indented);
% the author's signature in large calligraphy at the lower left with «عفي عنه» written small beside it.
\mspage{84}
\bab{\urdufont طريقۂ تعليم}
\begin{urduprose}
\pstart
مدرّسينِ كرام سے گذارش ہے كہ اس كتاب كے لئے فصل كے ساتھ دو گھنٹہ متعين فرمائيں پہلے گھنٹے ميں سرسرى انداز سے ضرورى ترجمہ كے ساتھ پڑھائيں اور اہم حصے ياد كرائيں دوسرے گھنٹے ميں كچھ تفصيل كے ساتھ سمجھائيں اور ہر سبق ميں اس غرض كو طلبہ كے ذہن نشيں كرائيں جس كے لئے سبق لكھا گيا اس سلسلے ميں تفہيم كے جو طريقے مناسب ہوں اختيار كريں مثلاً بورڈ \unclear{ميں} لكھنا بار بار دُہرانا وغيرہ۔ اس طريقے سے سال كے نصف اول ميں حصۂ صرف اور نصف آخر ميں حصۂ نحو پڑھائيں اور جب نحو شروع كرنے كا وقت آئے تو تين گھنٹے اس كے لئے فارغ فرمائيں۔ دو گھنٹوں ميں حسبِ سابق تعليم كا سلسلہ جارى ركھيں اور تيسرے گھنٹے ميں صرف و نحو كے آموختہ اسباق كى مشق كرائيں اور سنيں اس طرح بچّوں كو صرف و نحو ميں انشاء الله مہارت پيدا ہوگى۔
\pend
\pstart\raggedleft
محمد ميران عفي عنه
\pend
\end{urduprose}
